\section{Conclusion and Future Work}
\label{sec:conclusion}

This paper addressed the core systems problem behind multi-cloud FinOps decision support: provider billing exports are not directly comparable, commitment semantics distort raw workload cost, and actionable optimization requires accountable attribution rather than descriptive spend visibility alone. The implemented system responds with a canonical Unified Cost Allocation Schema, commitment-aware Net Effective Cost modeling, explicit shared-cost allocation semantics, and a deterministic decision pipeline that materializes waste signals, ranked recommendations, month-end forecasts, and lifecycle-ready validation outputs.

The current repository state provides meaningful internal evidence for this design. The synthetic benchmark materializes 18{,}992 normalized billing rows across AWS, Azure, and GCP, preserves Azure shared-cost fan-out exactly in each observed month, produces 95 persisted recommendations, and achieves forecast backtest quality of \$7.86 mean absolute error, 2.98\% mean absolute percentage error, and a 95.59\% within-10-percent rate. These results support the claim that the system is reproducible, auditable, and sufficiently coherent for workshop or conference-style systems evaluation.

The paper should nevertheless state its boundaries clearly. The benchmark is synthetic rather than enterprise production data. Idle compute detection is based on billing-side proxies rather than runtime telemetry. The lifecycle layer is implemented, but the current benchmark state is still pre-realization: all recommendations remain in the \texttt{recommended} state, so realized-savings validation is not yet mature. The system also stops at decision support and validation instrumentation; it does not perform autonomous remediation.

These limitations lead directly to future work. The most credible next step is enterprise validation on real multi-cloud billing exports with controlled access to ground-truth financial review. A second direction is telemetry-backed waste detection that combines billing semantics with utilization signals from compute, storage, and orchestration layers. A third is stronger closed-loop calibration, including richer implemented and verified lifecycle outcomes for recommendation classes such as \texttt{resize\_down} and \texttt{remove\_resource}. A fourth is workflow integration with ticketing, approval, and governance systems so that recommendations can move beyond local persistence into operational execution. Finally, the deterministic methods here provide a useful baseline for studying when more advanced forecasting or reasoning models improve accuracy without sacrificing auditability.
